\section{The marked family and a middle-dimensional class}
\label{hodge:section}

We first put all split Weil-type eightfolds for a fixed imaginary quadratic
field into one rational model.  The same model supplies the homology class
and the space of two-forms used in the symplectic construction.

\subsection{A model for every split Hermitian space}

Fix a positive squarefree integer $d$ and put $K=\Q(\sqrt{-d})$.
Choose an integral quaternion of norm $d$, which exists by the four-squares
theorem.  Its left-multiplication matrix $Q$ satisfies
$Q^tQ=QQ^t=d1_4$.  On $F=\Q^8$ set
\begin{equation}
 D=\begin{pmatrix}0&-Q^t\\Q&0\end{pmatrix},
 \qquad S_-=\diag(1_4,-1_4).
 \label{hodge:matrices}
\end{equation}
Thus $D^t=-D$, $D^2=-d1_8$, and $DS_-=-S_-D$.
Let
\[
 H=F_q\oplus F_p=F\otimes S,
 \qquad S=\Q e\oplus\Q f,
 \qquad \psi_0=dq^t\wedge dp.
\]
The action of $\sqrt{-d}$ is $D$ on both copies of $F$.  Our convention is
$\psi_0(x,Ix)>0$ for a positive compatible complex structure $I$.
For any matrix $M$ we write
$dq^tM\wedge dp=\sum_{a,b}M_{ab}\,dq_a\wedge dp_b$.

Consider the marked family
\begin{equation}
 \begin{split}
 \mathcal U&=\{\Pi\in\operatorname{Mat}_8(\C):
       \Pi=\Pi^t,\ \Im\Pi>0,\ [\Pi,D]=0\},\\
 A_\Pi&=\C^8/(\Z^8+\Pi\Z^8),\qquad z=q+\Pi p.
 \end{split}
 \label{hodge:period-domain}
\end{equation}
The principal polarization has class $\theta$, identified with $\psi_0$
in the fixed torus marking.  The domain $\mathcal U$ is a nonempty convex
open subset of the complex vector space of symmetric matrices commuting
with $D$: it contains $i1_8$.

\begin{proposition}[Reduction to the marked family]
\label{hodge:isogeny-reduction}
Every split Weil-type polarized abelian eightfold with $K$-action is
$K$-compatibly isogenous to a member of \eqref{hodge:period-domain}.
Consequently, algebraicity of the full rational Weil space on every
$A_\Pi$ implies the assertion for every compatible polarization type.
\end{proposition}

\begin{proof}
First, \eqref{hodge:period-domain} describes exactly the positive compatible
complex structures on $(H_\R,\psi_0)$ commuting with $D$.  Indeed the real
Lagrangian $F_q$ is transverse to $I F_q$: if $Iy\in F_q$ with
$y\in F_q$, then $\psi_0(y,Iy)=0$, so $y=0$.  A basis of $F_q$ is
therefore a complex basis.  In that basis the other lattice generators
give a symmetric period matrix with positive imaginary part, by the
Riemann bilinear relations.  The action of $D$ on the first basis is its
fixed matrix $D$, so complex linearity on the second basis says
$D\Pi=\Pi D$.  The converse follows from the same relations.

The $K$-Hermitian form
\[
 h_0(x,y)=\psi_0(Dx,y)+\sqrt{-d}\,\psi_0(x,y)
\]
is nonsingular and has the $K$-isotropic half-space $F_q$.  At $\Pi=i1_8$
the two eigenvalues of $D$ on $H^{1,0}$ have multiplicities four and four.
Those multiplicities are locally constant and $\mathcal U$ is connected,
so every member has the required Weil type.

For an arbitrary tuple in the theorem, write $(V,h)$ for its rational
$K$-Hermitian space.  The prescribed isotropic $K$-subspace has dimension
four.  Choose a basis $e_1,\ldots,e_4$ of it and vectors
$f_1,\ldots,f_4$ with $h(e_i,f_j)=\delta_{ij}$.
If $B_{ij}=h(f_i,f_j)$, replace
\[
 f_i\quad\hbox{by}\quad
 f_i-\frac12\sum_j B_{ij}e_j.
\]
Using that $h$ is linear in its first variable and
$B_{ij}=\overline{B_{ji}}$, the new $f_i$ span an isotropic complement.
Thus $(V,h)$ is hyperbolic.  Applying this construction to $h_0$ gives a
$K$-linear isometry $(H,h_0)\simeq(V,h)$.  Taking the coefficient of
$\sqrt{-d}$ shows that it also identifies the alternating forms.
The positive complex structure on $V_\R$ therefore pulls back to a
point of $\mathcal U$.

The transported integral homology lattice and $\Z^{16}$ are full lattices
in the same rational space, hence commensurable.  Multiplying the rational
identification by a suitable positive integer yields a $K$-compatible
isogeny of the corresponding complex tori.  Both are abelian varieties,
so the isogeny is algebraic.  If $f:B\to A$ is such an isogeny, then
$f^*:H^*(A,\Q)\to H^*(B,\Q)$ is an isomorphism, its inverse is
$f_*/\deg f$, and both maps preserve algebraic rational cycle classes.
They identify the two determinant lines defining the Weil spaces because
$f$ commutes with $K$.  This proves the reduction without a restriction
on integral polarization type.
\end{proof}

\subsection{Hodge classes at a general period}

We regard all rational cohomology groups as subspaces of
$\bigwedge^*H^*$ by the marking.  In particular the Weil space is the
same rational two-plane $W_K$ throughout $\mathcal U$.

\begin{proposition}
\label{hodge:generic-hodge}
There is a countable collection of proper closed analytic subsets of
$\mathcal U$ outside which
\begin{equation}
 H^{2j}(A_\Pi,\Q)\cap H^{j,j}(A_\Pi)=
 \begin{cases}
  \Q\theta^j,&j\ne4,\\
  \Q\theta^4\oplus W_K,&j=4.
 \end{cases}
 \label{hodge:generic-classes}
\end{equation}
Here $0\le j\le8$.
\end{proposition}

\begin{proof}
For each rational vector $v\in\bigwedge^{2j}H^*$, the condition
$v\in F^j_{\mathrm{Hdg}}$ is closed analytic, since the Hodge filtration
is a holomorphic subbundle of the constant cohomology bundle.
As $v$ is rational, this is equivalent to type $(j,j)$.
There are only countably many such vectors.  It remains to identify those
which have type $(j,j)$ at every period.

The centralizer of $K$ in $\Sp(H_\R,\psi_0)$ is the unitary group of
$h_0$, of signature $(4,4)$.  This follows directly from the formula for
$h_0$ and its recovery of $\psi_0$.  It acts on the compatible complex
structures in the family.  Their Hodge circles lie in $\SU(4,4)$:
on the $K$-complex representation their two circle weights each occur
four times, so their determinants are one.  A circle is noncentral,
since both weights occur.  The Lie algebra generated by all its
$\SU(4,4)$-conjugates is a nonzero ideal in the simple real Lie algebra
$\mathfrak{su}(4,4)$, hence is the whole algebra.  Thus a class fixed
by every Hodge circle is fixed by the connected group $\SU(4,4)$.

After complexification this is the $\SL_8$-representation on a standard
space $E$ and its dual.  Its exterior algebra decomposes as
\[
 \bigwedge^k(E\oplus E^*)=
 \bigoplus_{a+b=k}\bigwedge^aE\otimes\bigwedge^bE^*.
\]
An invariant in the $(a,b)$ summand is an intertwiner
$\bigwedge^bE\to\bigwedge^aE$.  The exterior powers are irreducible
and mutually inequivalent for $1\le a\le7$, whereas the zeroth and
eighth powers are both trivial.  The resulting invariant lines are the
matched summands $a=b$, together with $(a,b)=(8,0),(0,8)$.
The matched invariants are the powers of the nondegenerate pairing,
which is $\theta$ up to a nonzero scalar.  The two remaining lines are
exactly $W_K\otimes\C$.  Conversely $\theta$ and these determinant
lines have the asserted Hodge types at every period.  Descent to rational
cohomology gives the common spaces in \eqref{hodge:generic-classes}.
Exclude the proper Hodge locus of each remaining rational vector.
\end{proof}

\subsection{The symplectic model and its exceptional class}

The construction takes place on the separate symplectic torus
\begin{equation}
 \begin{split}
 T&=(\R^8_q/\Z^8)\times(\R^8_p/\Z^8),\qquad
       \omega=dq^t\wedge dp,\\
 \Omega&=\{dq^tM\wedge dp:M\in\operatorname{Mat}_8(\Q),\
                    M=M^t,\ [D,M]=0\}
       \subset\bigwedge^2H^*.
 \end{split}
 \label{hodge:omega-space}
\end{equation}
Use its symplectic orientation, the constant complex volume for $q+ip$,
and spin background zero.  Fix an integer $h_0\ge3$ and orient the graphs
\begin{equation}
 G_k=\{p=h_0kq\},\quad k\ge1,\qquad C=\{p=S_-q\}
 \label{hodge:graph-classes}
\end{equation}
by their $q$-coordinates.  Write $g_k,c\in\bigwedge^8H=H_8(T,\Q)$
for their fundamental classes.  Their slope matrices are symmetric, so
the graphs are Lagrangian.

Let $F_\C=U_+\oplus U_-$ be the eigenspace decomposition of $D$, with
eigenvalues $\pm i\sqrt d$, and set $H_\pm=U_\pm\otimes S_\C$.
Each $U_\pm$ has dimension four and each $H_\pm$ has dimension eight.

\begin{lemma}
\label{hodge:exceptional-homology}
There exists a rational class $\alpha_{\mathrm{ex}}\in\bigwedge^8H$
whose complexification belongs to
$\bigwedge^8H_+\oplus\bigwedge^8H_-$ and satisfies
\begin{equation}
 \alpha_{\mathrm{ex}}\wedge c\ne0,
 \qquad \alpha_{\mathrm{ex}}\wedge g_k=0\quad(k\ge1).
 \label{hodge:exceptional-pairings}
\end{equation}
Both of its pure eigenspace components are nonzero.
\end{lemma}

\begin{proof}
The sum of the two determinant lines descends to a rational two-plane:
the eigenspace splitting is defined over $K$, and its Galois involution
exchanges the two lines.  Projection of the graph vector space of $C$
onto $H_-$ is an isomorphism.  Indeed, writing $x=x_++x_-$, it is
\[
 x\longmapsto x_-\otimes e+(S_-x_+)\otimes f,
\]
and $S_-:U_+\to U_-$ is invertible.  Thus wedging with $c$ is nonzero
on $\bigwedge^8H_+$ and hence on the rational two-plane.  Choose a rational
vector where it is nonzero.  Its two components are conjugate, so neither
vanishes.  Finally $g_k$ has eigenspace bidegree $(4,4)$, whereas each
component of $\alpha_{\mathrm{ex}}$ is pure of degree eight; their wedge
is zero by dimension.
\end{proof}

Set
\begin{equation}
 \alpha=\alpha_{\mathrm{ex}}+g_1+g_2+g_3.
 \label{hodge:alpha}
\end{equation}
The three scalar slopes have a specific purpose: they detect the three
coefficients of a binary quadratic form.  This gives the two identities
which control the topology of the eventual Lagrangian.

\begin{lemma}
\label{hodge:linear-algebra}
The space $\Omega$ has dimension sixteen, and
\begin{equation}
 \ker\left(\bigwedge^2H^*\longrightarrow\bigwedge^6H,
                \ \eta\longmapsto\iota_\eta\alpha\right)=\Omega,
 \qquad
 H^*\otimes\Omega\xrightarrow{\ \wedge\ }\bigwedge^3H^*
 \text{ is injective}.
 \label{hodge:linear-identities}
\end{equation}
\end{lemma}

\begin{proof}
We may extend scalars to $\C$.  Skewness of $D$ makes $U_+$ and $U_-$
isotropic for the dot product and pairs them perfectly.  A symmetric
matrix commuting with $D$ therefore determines, and is determined by,
an arbitrary pairing between $U_+$ and $U_-$.  In the two $S$-coordinates
the corresponding two-form has alternating $S$-factor.  Consequently
\begin{equation}
 \Omega_\C=
 U_+^*\otimes U_-^*\otimes\bigwedge^2S_\C^*
       \subset\bigwedge^2H_\C^*,
 \label{hodge:mixed-forms}
\end{equation}
which has dimension sixteen.  These forms contract to zero against the
pure determinant lines and against every scalar graph volume, so
$\Omega_\C$ lies in the stated kernel.

Conversely, suppose $\iota_\eta\alpha=0$.  The components of its value
in eigenspace bidegrees $(6,0)$ and $(0,6)$ come only from contracting
the two nonzero pure volumes of $\alpha_{\mathrm{ex}}$ by the pure parts
of $\eta$.  Contraction into a nonzero top exterior power is an
isomorphism from two-forms to the complementary exterior power.  Hence
both pure parts of $\eta$ vanish, and $\eta$ is mixed.

Put $s_j=e+h_0jf$ for $j=1,2,3$.
Project the $(3,3)$ component of the contraction to
\[
 \bigwedge^3U_+\otimes\bigwedge^3U_-\otimes
              \operatorname{Sym}^6S_\C
\]
by alternating in each $U$-factor and multiplying the six $S$-factors
symmetrically.  Contraction of $g_j$ projects, up to one fixed nonzero
factor and the fixed volume isomorphisms, to
\[
 \eta(-\otimes s_j,-\otimes s_j)\otimes s_j^6.
\]
The vectors $s_1^6,s_2^6,s_3^6$ are independent: their coefficients of
$e^6,e^5f,e^4f^2$ give a nonzero Vandermonde determinant.
Thus all three pairing-valued evaluations vanish.  Every coefficient
of $\eta(-\otimes s,-\otimes s)$ is a binary quadratic in $s$,
and a binary quadratic vanishing on these three distinct lines is zero.
The symmetric $S$-part of $\eta$ therefore vanishes, leaving exactly
\eqref{hodge:mixed-forms}.  This proves the kernel assertion.

For wedge injectivity, use bases $u_i$ of $U_+^*$ and $v_j$ of $U_-^*$.
The subscripts $e,f$ denote their four copies in $H_\C^*$, and a basis
of $\Omega_\C$ is
\[
 z_{ij}=u_{i,e}\wedge v_{j,f}-u_{i,f}\wedge v_{j,e}.
\]
The following projections of a wedge product detect the four possible
types of first factor independently:
\[
\begin{array}{c|c}
 \text{first factor}&\text{three-letter block}\\ \hline
 U_{+,e}^*&U_{+,e}^*\wedge U_{+,f}^*\wedge U_{-,e}^*\\
 U_{+,f}^*&U_{+,e}^*\wedge U_{+,f}^*\wedge U_{-,f}^*\\
 U_{-,e}^*&U_{+,e}^*\wedge U_{-,e}^*\wedge U_{-,f}^*\\
 U_{-,f}^*&U_{+,f}^*\wedge U_{-,e}^*\wedge U_{-,f}^*
\end{array}
\]
For example, the first projection of $u_{a,e}\wedge z_{ij}$ is
$-u_{a,e}\wedge u_{i,f}\wedge v_{j,e}$, and no other first-factor
type contributes to that block.  The three indices freely recover
every coefficient.  The other rows give the same argument.  Hence all
coefficients of an element in the wedge kernel vanish.
\end{proof}
